\documentclass[11pt]{article}

\usepackage[utf8]{inputenc}
\usepackage[T1]{fontenc}
\usepackage[french]{babel}
\usepackage{amsmath}
\usepackage{amssymb}
\usepackage{geometry}
\usepackage{booktabs}
\usepackage{hyperref}

\geometry{margin=2.5cm}


\begin{document}

\section{Format de sortie des données}

On considère les variables suivantes :

\[
A_{i,j} =
\begin{cases}
1 & \text{si l'étudiant } i \text{ obtient le cours } j \\
0 & \text{sinon}
\end{cases}
\]

où :

\begin{itemize}
\item $i \in \{1,\dots,n\}$ représente un étudiant
\item $j \in \{1,\dots,m\}$ représente un cours
\item $r_{i,j}$ est le rang du cours $j$ pour l'étudiant $i$
\end{itemize}

\section{Liste des objectifs}

\subsection{Minimiser la somme des rangs}

\[
\min \sum_{i=1}^{n} \sum_{j=1}^{m} r_{i,j} A_{i,j}
\]

\textbf{Avantages :}

\begin{itemize}
\item maximise la satisfaction moyenne globale des étudiants
\item certains étudiants peuvent recevoir uniquement des cours très bien classés
\end{itemize}

\textbf{Limites :}

\begin{itemize}
\item peut générer des inégalités importantes entre étudiants (non équitable)
\item certains étudiants peuvent recevoir uniquement des cours mal classés
\end{itemize}

\subsection{Maximiser le nombre d'étudiants recevant leur premier choix}

\[
\max \sum_{i=1}^{n} \sum_{j=1}^{m} \mathbf{1}(r_{i,j} = 1)\, A_{i,j}
\]

où $\mathbf{1}(r_{i,j} = 1)$ vaut $1$ si le cours $j$ est le premier choix de l'étudiant $i$, et $0$ sinon.

\textbf{Avantages :}

\begin{itemize}
\item permet d'améliorer fortement la satisfaction pour une partie des étudiants
\end{itemize}

\textbf{Limites :}

\begin{itemize}
\item peut ignorer complètement les préférences restantes des étudiants
\item certains étudiants peuvent ne recevoir que des cours très mal classés
\end{itemize}

\subsection{Minimiser la somme des pires rangs par étudiant}

Soit $w_i$ le pire rang attribué à l'étudiant $i$.

\[
\min \sum_{i=1}^{n} w_i
\]

\textbf{Avantages :}

\begin{itemize}
\item réduit les mauvaises affectations pour tous les étudiants
\item plus équilibré que le pire rang global
\end{itemize}

\textbf{Limites :}

\begin{itemize}
\item ne garantit pas une satisfaction globale optimale
\item nécessite d'introduire une variable supplémentaire
\end{itemize}

\subsection{Maximiser le nombre de rangs faibles}

\[
\max \sum_{i=1}^{n}\sum_{j=1}^{m} \mathbf{1}(r_{i,j} \leq T) A_{i,j}
\]

où $T$ est un seuil correspondant à un bon classement.

\textbf{Avantages :}

\begin{itemize}
\item favorise l'attribution de cours bien classés
\end{itemize}

\textbf{Limites :}

\begin{itemize}
\item dépend du choix du seuil
\item ignore les différences entre les rangs
\end{itemize}

\subsection{Minimiser le nombre de rangs élevés}

\[
\min \sum_{i=1}^{n}\sum_{j=1}^{m} \mathbf{1}(r_{i,j} \geq T) A_{i,j}
\]

où $T$ est un seuil correspondant à un mauvais classement.

\textbf{Avantages :}

\begin{itemize}
\item réduit les affectations très peu satisfaisantes
\end{itemize}

\textbf{Limites :}

\begin{itemize}
\item dépend du choix du seuil
\item ignore les différences entre les rangs
\end{itemize}

\subsection{Minimiser l'écart de la somme des rangs entre étudiants}

On définit la somme des rangs des cours qui sont attribués à un étudiant $i$:

\[
s_i = \sum_{j=1}^{m} r_{i,j} A_{i,j}
\]

L'objectif consiste alors à minimiser l'écart entre $S_{\max}$ et $S_{\min}$ respectivement la somme maximal et minimal parmi les étudiants :

\[
\min (S_{\max} - S_{\min})
\]

\textbf{Avantages :}

\begin{itemize}
\item favorise une répartition plus équitable des affectations
\item limite les situations où certains étudiants sont très désavantagés
\end{itemize}

\textbf{Limites :}

\begin{itemize}
\item peut réduire la satisfaction moyenne globale
\item nécessite l'introduction de variables supplémentaires
\end{itemize}

\subsection{Maximiser la satisfaction croissante sur les cours voulus}

On associe à chaque rang $r_{i,j}$ une valeur de satisfaction décroissante
avec le rang :

\[
p(r_{i,j}) = \frac{1}{r_{i,j}}
\]

Ainsi, les cours les mieux classés apportent une satisfaction plus forte
que les cours moins bien classés.

L'objectif consiste alors à maximiser la somme des satisfactions obtenues :

\[
\max \sum_{i=1}^{n} \sum_{j=1}^{m} p(r_{i,j}) A_{i,j}
\]

\textbf{Avantages :}

\begin{itemize}
\item valorise fortement l'obtention des cours les plus souhaités
\item distingue mieux les très bons rangs que la simple somme des rangs
\end{itemize}

\textbf{Limites :}

\begin{itemize}
\item le choix de la fonction $p$ reste arbitraire
\item ne pénalise pas directement les cours très mal classés
\end{itemize}

\subsection{Minimiser la pénalité croissante des cours non voulus}

On associe à chaque rang $r_{i,j}$ une pénalité croissante avec le rang :

\[
q(r_{i,j}) = r_{i,j}^2
\]

Ainsi, les cours mal classés sont beaucoup plus pénalisés que les cours
moyennement classés.

L'objectif consiste alors à minimiser la somme des pénalités :

\[
\min \sum_{i=1}^{n} \sum_{j=1}^{m} q(r_{i,j}) A_{i,j}
\]

\textbf{Avantages :}

\begin{itemize}
\item pénalise fortement les affectations très peu souhaitées
\item réduit les situations où certains étudiants reçoivent uniquement des cours mal classés
\end{itemize}

\textbf{Limites :}

\begin{itemize}
\item le choix de la fonction $q$ reste arbitraire
\item peut réduire la satisfaction moyenne globale si l'objectif est utilisé seul
\end{itemize}

\subsection{Maximiser une satisfaction pondérée avec pénalité croissante}

On associe à chaque rang $r_{i,j}$ :

\begin{itemize}
\item une valeur de satisfaction décroissante $p(r_{i,j})$ pour favoriser les cours bien classés
\item une pénalité croissante $q(r_{i,j})$ pour pénaliser les cours mal classés
\end{itemize}

Par exemple :

\[
p(r_{i,j}) = \frac{1}{r_{i,j}}
\qquad
q(r_{i,j}) = r_{i,j}^2
\]

L'objectif consiste alors à maximiser un score combiné :

\[
\max \sum_{i=1}^{n} \sum_{j=1}^{m}
\left( \alpha \, p(r_{i,j}) - \beta \, q(r_{i,j}) \right) A_{i,j}
\]

où $\alpha$ et $\beta$ sont deux coefficients permettant d'équilibrer
l'importance de la satisfaction et de la pénalité.

\textbf{Avantages :}

\begin{itemize}
\item combine la recherche de bons choix et l'évitement des très mauvais choix
\item permet de moduler finement le comportement du modèle
\end{itemize}

\textbf{Limites :}

\begin{itemize}
\item nécessite de choisir les fonctions $p$ et $q$
\item nécessite de régler les coefficients $\alpha$ et $\beta$
\end{itemize}

\subsection{Maximiser le nombre d'étudiants satisfaits}

On définit pour chaque étudiant $i$ le nombre de cours bien classés
qui lui sont attribués :

\[
g_i = \sum_{j=1}^{m} \mathbf{1}(r_{i,j} \leq T_{bon}) A_{i,j}
\]

et le nombre de cours mal classés :

\[
b_i = \sum_{j=1}^{m} \mathbf{1}(r_{i,j} \geq T_{mauvais}) A_{i,j}
\]

Un étudiant est considéré comme satisfait si $K \geq g_i$
et $L \leq b_i$, où :
\begin{itemize}
\item K : nombre minimum de cours qu'il voulait
\item L : nombre maximum de cours qu'il ne voulait pas
\end{itemize}

\[
z_i =
\begin{cases}
1 & \text{si l'étudiant } i \text{ est satisfait} \\
0 & \text{sinon}
\end{cases}
\]

L'objectif consiste à maximiser le nombre d'étudiants satisfaits :

\[
\max \sum_{i=1}^{n} z_i
\]

\textbf{Avantages :}

\begin{itemize}
\item favorise des affectations globalement acceptables pour un grand nombre d'étudiants
\end{itemize}

\textbf{Limites :}

\begin{itemize}
\item dépend du choix des seuils $T_{bon}$, $T_{mauvais}$, $K$ et $L$
\item ne distingue pas les niveaux de satisfaction au-delà de ces seuils
\end{itemize}

\section{Sélection de l'objectif}

Les objectifs présentés précédemment seront évalués entre 1 et 10 par les membres du projet selon plusieurs critères :

\begin{itemize}

\item \textbf{Compréhension pour un membre du projet}  
Le principe de l'objectif est-il clair pour quelqu'un qui travaille sur le projet ?

\item \textbf{Compréhension pour un non-membre du projet}  
Une personne extérieure (étudiant ou enseignant) peut-elle comprendre facilement l'objectif ?

\item \textbf{Pertinence pour représenter l'équité}  
L'objectif permet-il de répartir les cours de manière équitable entre les étudiants ?

\item \textbf{Facilité d'implémentation}  
L'objectif est-il simple à implémenter dans un solveur ?

\item \textbf{Envie d'utiliser cet objectif}  
Globalement, cet objectif semble-t-il pertinent et intéressant à utiliser ?

\end{itemize}

\end{document}